\documentclass[aspectratio=169,12pt]{beamer}
\usepackage{fontspec}
\usepackage[portuguese]{babel}
\usepackage{xcolor,listings,tikz,booktabs,amsmath,amssymb,graphicx,multicol,array,colortbl}
\usetikzlibrary{arrows.meta,positioning,shapes.geometric,calc}

\definecolor{navy}{HTML}{0D1B3E}
\definecolor{gold}{HTML}{F5A623}
\definecolor{qgreen}{HTML}{1E8B4C}
\definecolor{qblue}{HTML}{1A6EAE}
\definecolor{codebg}{HTML}{EEF3F8}

\usetheme{default}
\useinnertheme{circles}
\setbeamertemplate{navigation symbols}{}
\setbeamertemplate{footline}{%
  \hfill{\color{gray!60}\scriptsize\insertframenumber/\inserttotalframenumber}\hspace{.6em}\vspace{.4em}}

\setbeamercolor{frametitle}{bg=navy,fg=white}
\setbeamercolor{structure}{fg=navy}
\setbeamercolor{alerted text}{fg=gold}
\setbeamercolor{item}{fg=gold}
\setbeamercolor{subitem}{fg=navy}
\setbeamercolor{block title}{bg=qblue,fg=white}
\setbeamercolor{block body}{bg=qblue!8,fg=black}
\setbeamercolor{block title alerted}{bg=gold,fg=navy}
\setbeamercolor{block body alerted}{bg=gold!15,fg=navy}
\setbeamercolor{block title example}{bg=qgreen,fg=white}
\setbeamercolor{block body example}{bg=qgreen!10,fg=black}
\setbeamerfont{frametitle}{size=\large,series=\bfseries}
\setbeamerfont{block title}{series=\bfseries}

\setbeamertemplate{title page}{%
  \begin{tikzpicture}[remember picture,overlay]
    \fill[navy](current page.north west)rectangle(current page.south east);
    \fill[gold]([yshift=.7cm]current page.south west)rectangle(current page.south east);
    \node[anchor=south,yshift=.73cm,navy,font=\tiny\bfseries]at(current page.south)
      {INTENSIVÃO QUANT AI \textbullet{} IMPACT UFSCAR};
  \end{tikzpicture}
  \vspace{.65cm}
  \begin{center}
    {\color{gold}\scriptsize\bfseries INTENSIVÃO QUANT AI \textbullet{} IMPACT UFSCAR}\\[.4cm]
    {\color{white}\LARGE\bfseries\inserttitle}\\[.25cm]
    {\color{gold!85}\normalsize\insertsubtitle}\\[.5cm]
    {\color{white!70}\small\insertauthor}\\[.1cm]
    {\color{white!50}\footnotesize\insertdate}
  \end{center}}

\lstdefinestyle{python}{language=Python,
  basicstyle=\ttfamily\scriptsize\color{qblue},
  keywordstyle=\color{navy}\bfseries,
  stringstyle=\color{qgreen!80!black},
  commentstyle=\color{gray!70}\itshape,
  backgroundcolor=\color{codebg},
  frame=l,framerule=2pt,rulecolor=\color{gold},
  xleftmargin=1em,breaklines=true,showstringspaces=false,tabsize=4,
  extendedchars=false,inputencoding=ascii}
\lstset{style=python}

\author{Felipe Maldonado\\{\scriptsize VP Diretoria Quant~\textbullet{}~Impact UFSCAR}}
\date{Intensivão Quant AI 2025}
\title{Aula 10 --- Relatório e Defesa}
\subtitle{Relatório Final, Análise de Sensibilidade e Simulação de Defesa}

\begin{document}

\begin{frame}[plain]
\titlepage
\end{frame}

\begin{frame}[fragile]{Nesta Aula}
\begin{columns}[T]
  \column{0.5\textwidth}
\begin{block}{Teoria (20 min)}
\begin{itemize}
  \item Estrutura do relatório técnico para o Desafio Itaú
  \item Quatro dimensões avaliadas pelos juízes
  \item Como defender sob pressão: perguntas difíceis
  \item Erros comuns a evitar
\end{itemize}
\end{block}
  \column{0.47\textwidth}
\begin{block}{Código (80 min)}
\begin{itemize}
  \item \texttt{painel\_relatorio\_final()} — figura principal com 4 subplots
  \item \texttt{analise\_sensibilidade()} — lookback e n\_long\_pct
  \item \texttt{ic\_temporal()} — IC mensal e rolling 12m
  \item Montar todos os resultados em estrutura de relatório
  \item Simulação de defesa: perguntas e respostas
\end{itemize}
\end{block}
\end{columns}
\end{frame}

\begin{frame}[fragile]{Estrutura do Relatório Final}
\begin{columns}[T]
  \column{0.5\textwidth}
\begin{itemize}
  \item \textbf{1. Sumário Executivo} (1 pág.): hipótese, metodologia, resultado principal, limitações
  \item \textbf{2. Dados e Universo}: fonte, período, limpeza, survivorship bias
  \item \textbf{3. Metodologia}: sinal (fórmula + código), portfólio, backtest, rigor OOS
  \item \textbf{4. Resultados}: tabela IS/OOS, equity curve, IC, DSR
  \item \textbf{5. Sensibilidade}: variação de lookback e tamanho do portfólio
  \item \textbf{6. Limitações e Trabalhos Futuros}: honestidade sobre lacunas
  \item \textbf{7. Referências}: papers completos (APA ou ABNT)
\end{itemize}
  \column{0.47\textwidth}
\begin{block}{Critérios dos Juízes}
\begin{itemize}
  \item \textbf{Rigor metodológico}: walk-forward, shift(1), sem look-ahead
  \item \textbf{Fundamentação econômica}: por que momentum funciona?
  \item \textbf{Qualidade do código}: reprodutível, versionado no GitHub
  \item \textbf{Comunicação}: proporcional à evidência, limitações declaradas
\end{itemize}
\end{block}\vspace{.3em}
\begin{alertblock}{Regra de ouro}
Um resultado honesto com limitações bem articuladas vale mais que um resultado inflado com lacunas escondidas.
\end{alertblock}
\end{columns}
\end{frame}

\begin{frame}[fragile]{Pipeline Completo — Resumo}

\begin{center}
\begin{tikzpicture}[node distance=.55cm and .3cm,
  box/.style={draw=navy,fill=navy!8,rounded corners=3pt,
              text width=1.55cm,align=center,font=\tiny\bfseries,inner sep=3pt},
  lbl/.style={font=\tiny,gold},
  arr/.style={-Stealth,navy,thick}]
  \node[box](d){Dados\\yfinance};
  \node[box,right=of d](e){EDA\\análise};
  \node[box,right=of e](s1){Sinal v1\\12-1};
  \node[box,right=of s1](s2){Sinal v2\\vol-adj};
  \node[box,right=of s2](p){Portfólio\\top 20\%};
  \node[box,right=of p](b){Backtest\\IS};
  \node[box,right=of b](v){WF-OOS\\DSR};
  \node[box,right=of v](g){GenAI\\narrativa};
  \node[box,right=of g](r){Relatório\\defesa};

  \draw[arr](d)--(e);\draw[arr](e)--(s1);\draw[arr](s1)--(s2);
  \draw[arr](s2)--(p);\draw[arr](p)--(b);\draw[arr](b)--(v);
  \draw[arr](v)--(g);\draw[arr](g)--(r);

  \node[lbl,below=.3cm of d]{Aula 2};
  \node[lbl,below=.3cm of e]{Aula 3};
  \node[lbl,below=.3cm of s1]{Aula 4};
  \node[lbl,below=.3cm of s2]{Aula 7};
  \node[lbl,below=.3cm of p]{Aulas 5-6};
  \node[lbl,below=.3cm of b]{Aula 5};
  \node[lbl,below=.3cm of v]{Aula 8};
  \node[lbl,below=.3cm of g]{Aula 9};
  \node[lbl,below=.3cm of r]{Aula 10};
\end{tikzpicture}
\end{center}
\vspace{.2em}
\begin{columns}[T]
  \column{0.48\textwidth}
  \begin{block}{Parquets produzidos}
    \texttt{precos\_ibov} $\cdot$ \texttt{ret\_diarios\_limpo} $\cdot$ \texttt{ret\_mensais\_limpo}\\
    \texttt{sinal\_v1} $\cdot$ \texttt{sinal\_v2} $\cdot$ \texttt{pesos\_v1} $\cdot$ \texttt{pesos\_v2}\\
    \texttt{retorno\_carteira} $\cdot$ \texttt{retorno\_walkforward\_liquido}
  \end{block}
  \column{0.48\textwidth}
  \begin{exampleblock}{Entregáveis do Desafio}
    Relatório PDF $\cdot$ Código no GitHub\\
    Apresentação 10 min + 5 min Q\&A\\
    Notebook reprodutível do zero
  \end{exampleblock}
\end{columns}
\end{frame}

\begin{frame}[fragile]{Como Defender — Perguntas Difíceis}
\begin{columns}[T]
  \column{0.5\textwidth}
\begin{block}{Pergunta 1}

\textit{``Seu Sharpe é atribuível ao momentum ou ao beta de mercado?''}\\[.3em]
\textbf{Resposta}: Compare correlação da carteira com IBOVESPA. Mostre alpha de Jensen $\alpha = r_p - \beta\,r_m$.

\end{block}\vspace{.3em}
\begin{block}{Pergunta 2}

\textit{``Como controlaram survivorship bias?''}\\[.3em]
\textbf{Resposta}: Declarar como limitação. Explicar que yfinance com histórico minimiza, mas não elimina. Citar na seção 6.

\end{block}\vspace{.3em}
\begin{block}{Pergunta 3}

\textit{``Por que equal-weight e não Markowitz?''}\\[.3em]
\textbf{Resposta}: DeMiguel et al.\ (2009) — 1/N vence Markowitz OOS. Mostrar tabela de Sharpe comparativo.

\end{block}
  \column{0.47\textwidth}
\begin{block}{Regras da Defesa}
\begin{itemize}
  \item Nunca finja saber o que não sabe
  \item \textit{``Boa pergunta, não testamos, mas hipotetizamos que...''} é melhor que erro com confiança
  \item Cite papers para todo claim quantitativo
  \item Mostre os gráficos — visuais vencem palavras
  \item Declare limitações antes de te perguntarem
\end{itemize}
\end{block}\vspace{.3em}
\begin{exampleblock}{Encerramento}

Em 10 semanas: dados reais, sinal acadêmico, backtest rigoroso, validação OOS, GenAI.
\textbf{Façam bonito no Desafio!}
\end{exampleblock}
\end{columns}
\end{frame}

\begin{frame}[fragile]{Referências — Pipeline Completo}
\begin{multicols}{2}
\tiny
\begin{itemize}
  \item \textbf{Jegadeesh, N. \& Titman, S.} (1993). Returns to buying winners. \textit{Journal of Finance}, 48(1).
  \item \textbf{Moskowitz, T. et al.} (2012). Time series momentum. \textit{JFE}, 104(2), 228--250.
  \item \textbf{Fama, E. \& French, K.} (1993). Common risk factors. \textit{JFE}, 33(1), 3--56.
  \item \textbf{DeMiguel, V. et al.} (2009). Optimal versus naive diversification. \textit{RFS}, 22(5).
  \item \textbf{Bailey, D. \& López de Prado, M.} (2014). The deflated Sharpe ratio. \textit{JPM}, 40(5).
  \item \textbf{Harvey, C. et al.} (2016). Cross-section of expected returns. \textit{RFS}, 29(1).
  \item \textbf{López de Prado, M.} (2018). \textit{Advances in Financial Machine Learning}. Wiley.
  \item \textbf{Daniel, K. \& Moskowitz, T.} (2016). Momentum crashes. \textit{JFE}, 122(2).
  \item \textbf{Cont, R.} (2001). Empirical properties of asset returns. \textit{Quantitative Finance}, 1(2).
  \item \textbf{Markowitz, H.} (1952). Portfolio Selection. \textit{Journal of Finance}, 7(1).
\end{itemize}
\end{multicols}
\end{frame}


\end{document}
